%*******************************************************
% Abstract in English
%*******************************************************
\pdfbookmark[0]{Abstract}{Abstract}

\begin{otherlanguage}{american}
	\chapter*{Abstract}    

    % Inhalt
    % Fragestellung: welche Fragestellung(en) versucht der Text zu beantworten, was ist der Umfang der Forschung, was sind die zentralen Argumente und Behauptungen?
    This bachelor thesis investigates the suitability and significance of the presented dash\-board solution as a modernization measure to sustainably improve the usability of legacy monitoring processes. The term \textit{usability} is defined by DIN EN ISO 9241-110 and subdivided into the three dimensions of \textit{effectiveness}, \textit{efficiency}, and \textit{user satisfaction}. The basis of this study is the inefficient usability of the result view of a legacy test pool, which is nevertheless used for daily system status monitoring.

    % Methodik
    % Methodologie: welche Methoden/Zugänge nutzt der Autor/die Autorin, auf welche empirische Basis stützt sich der Text?
    To evaluate the suitability of the newly developed dash\-board, its impact on the improvement of the three usability dimensions was examined. For this purpose, a two-part user study was conducted and its results statistically evaluated. The first part of the study consisted of a user survey, whose questionnaire was based on Jochen Prümper’s survey to verify the implementation of the design recommendations of DIN EN ISO 9241-110. The second part comprised a task-based usability test derived from the survey, which in turn formed the basis for the users’ evaluation of the system.

    % Ergebnis-Teaser
    % zu welchen Ergebnissen kam die im Text dargestellte Forschungsarbeit, was sind die zentralen Schlussfolgerungen des Textes?
    Within the user study, the new interface, which compresses the test results of many views of the legacy system into just two, convinced with a clear and intuitive presentation. This led to shorter processing times and greatly improved user satisfaction. Although the dash\-board was the cause of an increased error rate in the usability test, an overall improvement of the monitoring process was evident. The solution is therefore classified as a suitable and non-invasive modernization method, but shall serve exclusively as an initial step toward system optimization within the modernization process.

    
%	A short summary of the contents in English of about one page. The following points should be addressed in particular:
%	\begin{itemize}
%		\item Motivation: Why did this work come about? Why is the topic of the work interesting (for the general public)? The motivation should be abstracted as far as possible from the specific tasks that may be given by a company.
%		\item Content: What is the content of this thesis? What exactly is covered in the thesis? The methodology and working method should be briefly discussed here.
%		\item Results: What are the results of this work? A brief overview of the most important results as a teaser to read the complete thesis.
%	\end{itemize}
%	\medskip
	
%	\noindent
%	BTW: A great guide by Kent Beck how to write good abstracts can be found here:
%	\begin{center}
%		\url{https://plg.uwaterloo.ca/~migod/research/beckOOPSLA.html}
%	\end{center}
\end{otherlanguage}
